\chapter{Mitigating microplastics vertical migration in usaturated soil with biochar}
\label{chap:five}

\section{Introduction} \label{sec:5.1}
In the previous chapters, it is demonstrated that biochar can effectively improve the water retention capacity of unsaturated soils. In this chapter, building upon this effect, biochar is further proposed as a functional material for controlling the migration of \ac{MPs} pollutants in soils subjected to \acp{WDC}. \par
Plant-based biochar demonstrates promising capability in adsorbing \ac{MPs} within soil porous media. However, the function of biochar in mitigating MPs’ vertical penetration during \acp{WDC}, which is typical of seasonal precipitation and evaporation, remains uncertain. Furthermore, few studies have investigated the mechanisms by which biochar interacts with soil. In this chapter, a series of column tests are conducted to assess the \ac{MPs}'s retention capabilities of soil-biochar porous media under saturated and wetting-drying conditions. The water retention and hydrophilic properties are investigated to elucidate the impact of wetting and drying cycles. Additionally, different biochar-soil structures are compared to optimise the structural design. The mechancism of biochar's effect on \ac{MPs} retention is further revealed by micro-morphological and surface chemical analyses via SEM and FTIR tests. The findings of this chapter provide insights into the application of biochar in mitigating \ac{MPs} pollution in unsaturated soil environments, which is crucial for protecting groundwater quality and ecological health. \par

\section{Material characterisation} \label{sec:5.2}
\subsection{Testing soil and biochar}
The soil tested in this study was collected near the UK's HS2 railway line, close to the University of Warwick. The soil was first sieved through a 2 mm mesh. Figure \ref{fig:PSDsandBiochar} shows the soil's particle size distribution (PSD), as measured according to the British Standard (BS 1377-2). The liquid limit and plastic limit of the soil are 24.26\% and 14.60\%, respectively. This soil exhibits low plasticity with a plasticity index of 9.66. The soil is classified as sandy loam (BS 1377-2). The organic matter content is 5.0\% and the nutrient content is 0.171\%. The soil performs alkalescence with a pH value of 7.8. To minimise the uncertain impacts caused by the chemical inhomogeneity at the soil particle surface, such as impurities, ogamic matters, microorganisms and oxides, the soil was first soaked in a concentrated HCl solution for 24 hours and then in a NaOH solution for 24 hours (Hsieh et al., 2022; Jiang et al., 2018; Tong et al., 2020). Then, it was flushed with deionised water until the pH reached around 7. Afterwards, the soil was oven-dried for 12 hours at 60 $^\circ\text{C}$ for preparation. \par
\begin{sloppypar}
The wood-derived biochar was purchased from SoilFixer Ltd, UK. The biochar was produced from arboricultural arisings (from mixed European hardwood species, such as oak and robinia) through a pyrolysis process (4-8 hours) in a retort kiln at 350-500 $^\circ\text{C}$. Schruff (2018) proposed that particle size significantly influences the migration behaviour of MP particles. Therefore, the biochar was pre-sieved using a 2.0 mm sieve and adjusted to have a similar particle size distribution (identical values of d30, d50 and d60) as the testing soil, as shown in Figure \ref{fig:PSDsandBiochar}. \par
\end{sloppypar}
\begin{figure}[htbp]
    \centering
    \includegraphics[width=0.65\linewidth]{../figures/Ch5/ch5Figures/PSDsandBiochar.jpg}
    \caption{Particle size distribution of sandy loam and biochar}
    \label{fig:PSDsandBiochar}
\end{figure}

\subsection{Soil and biochar mixture}
Table \ref{tab:physicalSB} summarises the liquid limit (LL), optimal moisture content (OMC) and specific gravity of soil-biochar mixture with different biochar contents. Besides, the morphological characteristics of the uncontaminated biochar and soil mixtures were examined using a Scanning Electron Microscope (SEM), as illustrated in Figure \ref{fig:SEMsandBiochar}. The sand particles, marked in yellow, predominantly exhibited rounded shapes. In contrast, some soil particles show slight wrinkles and rough surfaces. The biochar particles, traced in blue, displayed a variety of shapes, including chips, chunks and sheets. Biochar particles also exhibited an alveolate structure resembling a honeycomb (Bataillou et al., 2022), featuring open, interconnected chambers that enhance capillary action and water retention. Additionally, the large surface area of biochar significantly contribute to its capacity for adsorbing microplastic particles Mohan et al. (2014). \par
\begin{table}[htbp]
    \centering
    \caption{Physical properties of soil-biochar mixture}
    \begin{tabularx}{\textwidth}{
        c
        >{\justifying\centering\arraybackslash}m{3cm}
        >{\justifying\centering\arraybackslash}m{4cm}
        >{\justifying\centering\arraybackslash}m{2.5cm}
    }
    \toprule
    Biochar content (\%) 
        & Liquid limit (LL, \%) 
        & Optimal moisture content (by mass, \%) 
        & Specific gravity 
    \\ \midrule
    0  & 24.26 & 13.6 & 2.56 \\
    5  & 26.27 & 14.6 & 2.46 \\
    10 & 30.67 & 16.5 & 2.36 \\
    15 & 33.82 & 18.6 & 2.28 \\ \bottomrule
    \end{tabularx}
    \label{tab:physicalSB}
\end{table}
\begin{figure}[htbp]
    \centering
    \includegraphics[width=0.8\linewidth]{../figures/Ch5/ch5Figures/SEMsandBiochar.jpg}
    \caption{SEM images of soil and biochar mixture}
    \label{fig:SEMsandBiochar} % Fig.2
\end{figure}

\subsection{Microplastic sphere and suspension}
Polystyrene (PS) is one of the most prevalent microplastics in natural waters and is commonly used in commercial packaging and the automotive industry (Schirinzi et al., 2019) (Lv et al., 2024). This study selected monodispersed \qty{1}{\micro\meter} PS microspheres with a density of 1.05 g/cm³ to prepare the MPs suspension. The Fourier-transform infrared spectra (FTIR) test was performed to reveal the hydrophilic nature. The results (Figure \ref{fig:FTIRps}) indicate that the PS spheres lack significant hydrophilic groups, such as hydroxyl or carboxyl groups, confirming their hydrophobic nature. Therefore, surfactants were utilised to improve the dispersion of MPs in water (Ahmad et al., 2023). The suspension was then subjected to ultrasonic bathing to achieve a homogenous dispersion of MP particles. The concentration of MP suspension was 10 mg/L. To simulate the ionic strength of rainwater, the ionic strength of the MP suspension was adjusted using NaCl (0.84 mM) and NaHCO3 (0.16 mM) (O'Connor et al., 2019). \par
\begin{figure}[htbp]
    \centering
    \includegraphics[width=0.65\linewidth]{../figures/Ch5/ch5Figures/FTIRps.jpg}
    \caption{FTIR spectra of polystyrene sphere}
    \label{fig:FTIRps} % Fig.3
\end{figure}


\section{Testing methods} \label{sec:5.3}
\subsection{Sample preparation and testing conditions}
Samples were prepared in cylindrical polyvinyl chloride (PVC) columns with a diameter of 30 mm and a height of 225 mm. The soil and biochar were mixed at the optimal moisture contents and then carefully compacted in three layers to achieve predetermined initial porosities of 0.3, 0.4 and 0.5 following the equation:
\begin{equation}
     m=h\cdot G_{a}\cdot \rho_{w}\cdot S\cdot \left( 1-n_{0} \right)
    \label{eq:sampleWeigth}
\end{equation}
\begin{equation}
     G_{a}=\frac{\left ( 1-\xi  \right )\cdot G_{s}+\xi\cdot G_{b}}{100}
    \label{eq:specificGravity}
\end{equation}
where \textit{h} is the height of the biochar-soil sample (cm), \textit{m} is the total weight of the sample (g), $G_a$ is the specific gravity of soil-biochar mixture, $\rho_w$ is the density of water (g/cm$^3$), \textit{S} is the section area of the sample (cm$^2$), $n_0$ is the target initial porosity, $G_s$ is the specific gravity of soil particles, $G_b$ is the specific gravity of biochar and $\xi$ is the biochar content (\%). At the bottom of the sample, a 20 mm cushion layer of coarse sand was added to make the effluent flow out evenly. A fabric screen was stuck to the bottom of the PVC column to prevent large particles from escaping. After compaction, the samples were soaked in deionised water inside a vacuum container to achieve saturation. Four environmental and physical factors were designed: wetting-drying cycle, biochar content, initial porosity and biochar-soil mixing format. As shown in Figure \ref{fig:sample}, the mixing formats of biochar-soil refer to either a direct mixture or sandwiched layers. The testing details of 15 specimens are shown in Table \ref{tab:testCondition}.\par
\begin{figure}[htbp]
    \centering
    \includegraphics[width=0.8\linewidth]{../figures/Ch5/ch5Figures/sample.jpg}
    \caption{Configuration of the column test}
    \label{fig:sample} % Fig.4
\end{figure}

\begin{table}[htbp] % 这个表格要修改格式，定稿的时候要修改 因为他可能没有在相应的subsection出现
    \centering
    \caption{Testing details of column tests}
    \renewcommand{\arraystretch}{1.3}
    \begin{tabularx}{\textwidth}{
        >{\centering\arraybackslash}m{1.4cm}   % ID
        >{\centering\arraybackslash}m{2.4cm}   % mixing format
        >{\centering\arraybackslash}m{2.0cm}   % porosity
        >{\centering\arraybackslash}m{2.4cm}   % biochar content
        >{\centering\arraybackslash}m{4.5cm}   % Wetting-drying cycles
    }
    \toprule
    ID &
      Soil-biochar mixing format &
      Initial porosity &
      Biochar content (\%) &
      Wetting-drying cycles \\
    \midrule

    ML0 &
      \multirow{12}{*}{Mixed} &
      \multirow{4}{*}{0.50} &
      0 &
      \multirow{12}{*}{
        \parbox{3.8cm}{\justifying
        Each group was tested under constant saturated conditions and six wetting-drying cycles.}
    } \\
    ML05 &  &                       & 5  &  \\
    ML10 &  &                       & 10 &  \\
    ML15 &  &                       & 15 &  \\
    \cmidrule(lr){3-3}

    MM0  &  & \multirow{4}{*}{0.40} & 0  &  \\
    MM05 &  &                       & 5  &  \\
    MM10 &  &                       & 10 &  \\
    MM15 &  &                       & 15 &  \\
    \cmidrule(lr){3-3}

    MD0  &  & \multirow{4}{*}{0.30} & 0  &  \\
    MD05 &  &                       & 5  &  \\
    MD10 &  &                       & 10 &  \\
    MD15 &  &                       & 15 &  \\
    \cmidrule(lr){2-3} \cmidrule(l){5-5}

    LL05 &
      \multirow{3}{*}{Layered} &
      \multirow{3}{*}{0.50} &
      5 &
      \multirow{3}{*}{Constantly saturated.} \\
    LL10 &  &                       & 10 &  \\
    LL15 &  &                       & 15 &  \\
    \bottomrule
    \end{tabularx}
    \label{tab:testCondition}
\end{table}

\subsection{Soil water retention characteristics tests}
% 参考审稿回复意见，补充仪器照片和降水的数据来确定研究工况！！！！！！！！！！！！！！
In the compacted soil, pore water can be redistributed by filling in smaller pores due to the capillary effect (Mualem, 1976). Therefore, the pore size distribution, including existing cracks in soil samples, can be indicated by measuring the change in the soil water retention curves (Sun and Cui, 2018; Zhai et al., 2023). To determine the number of wetting-drying cycles used in the leaching tests, a series of soil water retention curve (SWRC) tests was conducted. The number of cycles at which the soil’s desorption curves tend to stabilise at a certain state was chosen, which reflects the propagation of cracks reaching a stable state accordingly. The water retention characteristics of soil-biochar samples were measured using a WP4C dewpoint potentiometer and a HYPROP moisture release system. WP4C was used for the matric suction range from 250 kPa to 100 MPa, while HYPROP could accurately measure matric suction less than 250 kPa. Azizi et al. (2021) have suggested that combining data from both instruments is reliable for assessing the relationship between the soil’s water content and its matric suction. During the test, samples were wetted by gradually adding water and dried in ovens.\par

\subsection{Leaching tests}
The leaching tests were conducted at constant laboratory conditions of 25 ℃ and 40\% relative humidity. Before leaching the MP suspension, the samples are pre-leached with the simulated rainwater of 10-pore volume to establish ionic equilibrium. Subsequently, the MP suspension was leached from the specimen and filtered through the sample under gravity. The simulated extreme rainfall intensity (RI) was 20 mm/h, which was reported as a critical RI that can potentially cause flash floods (Kendon et al., 2023). The rainfall intensity was controlled by adjusting the flow rate of the dosing pump, maintaining it at a targeted value. The effluent was collected at 10-minute intervals. For the samples designed to undergo wetting-drying cycles, after completing the leaching process (wetting), the samples were oven-dried (drying) and then subjected to the subsequent wetting process using MP suspension. Based on the results of SWRC tests, the SWRC remains at a constant state after five wetting-drying cycles. Lu et al. (2021) also demonstrated that the propagation of soil cracks tended to be stable after around five wetting-drying cycles. Based on our test results and the existing literature, this study performed six wetting-drying cycles, considering certain contingencies, during the leaching process. \par

\subsection{MP particles extraction and quantification}
To quantify the MP particles in effluent and the samples, an MPs-solvent density gradient method (see Figure \ref{fig:testProcedure}) was employed for MP extraction and counting (Ben et al., 2024). After the leaching test, each sample was equally segmented into five sections and oven-dried (40 ℃). A 3-4g piece was collected from each session. The sample piece was flushed with an alcohol solution (0.80g/cm³) to remove most impurities. The effluent was collected and centrifuged at 3000 g (g is the gravitational acceleration) for 30 minutes. During the process, MP particles with a density of 1.05 g/cm$^3$ settled and assembled at the bottom. The upper limpid liquid was carefully sucked away until 50 mL remained, and then 1 mL of cetyltrimethylammonium bromide surfactant was added. The suspension was placed in an ultrasonic water bath to promote the dispersion. Finally, \qty{20}{\micro\liter} of the solution was transferred to a hemacytometer. The MP spheres were counted through a microscope (O'Connor et al., 2019). The PS sphere and a few remaining impurities (biochar or soil particles) can be easily distinguished through their morphological difference. \par
To quantify MPs in effluents, the effluents of leaching tests were placed in a laboratory environment for natural evaporation (Ahmad et al., 2023). The remaining operations follow steps 6 and 7, as depicted in Figure \ref{fig:testProcedure}. Additionally, a session of dry samples was collected (step 1) to conduct an SEM test and investigate the retained MP particles. \par
\begin{figure}[htbp]
    \centering
    \includegraphics[width=1.0\linewidth]{../figures/Ch5/ch5Figures/testProcedure.jpg}
    \caption{Protocol for extracting and counting microplastics}
    \label{fig:testProcedure} % Fig.5
\end{figure}


\section{Results and analysis} \label{sec:5.9}
\subsection{Effect of wetting-drying cycle}
Leaching tests were conducted on two groups of samples in constant saturated or wetting-drying cycle conditions. Fig. 6 shows the breakthrough curves at an initial porosity of 0.5. These curves initially show an upward trend with increasing leaching suspension volume, eventually stabilising upon reaching adsorption-desorption equilibrium between the adsorbent and MP particles. With identical biochar content, the breakthrough curves under wetting-drying cycles exhibit higher values than those under saturated conditions, suggesting that more MP particles penetrated through the soil-biochar samples due to the wetting-drying cycles. When biochar content is increased from 0 to 15\%, the equilibrium values of the breakthrough curves decrease by around 50\%. Additionally, the breakthrough curves under wetting-drying cycles reach the equilibrium values slowly and exhibit slight fluctuations when the samples are dried and rewetted. Therefore, the adsorption on MPs significantly depends on hydraulic conditions and the biochar content. \par
\begin{figure}[htbp]
    \centering
    \includegraphics[width=1.0\linewidth]{../figures/Ch5/ch5Figures/BreakthroughCurve.jpg}
    \caption{Breakthrough curves of MPs suspension under constant saturated conditions and wetting-drying cycles with the initial porosity of (a) $n_0$= 0.5, (b) $n_0$= 0.4 and (c) $n_0$= 0.3}
    \label{fig:BreakthroughCurve} % Fig.6
\end{figure}
An integration method by Tong et al. (2020) is employed here to calculate the cumulative MPs proportion in the effluent (C$_{ex}$, \%). The additional MPs due to wetting-drying cycle are quantified by a parameter, $\Delta$ (\%), calculated by the following equations:
\begin{equation}
      C_{ex}=\frac{1}{N}\cdot \int_{0}^{N}\frac{C\left ( n \right )}{C_{0}}\textrm{d}n\cdot 100
    \label{eq:Cex}
\end{equation}
\begin{equation}
      \Delta = \int_{0}^{6}\frac{C_{\textbf{wd}}\left ( n \right )}{C_{0}}\textrm{d}n-\int_{0}^{6}\frac{C_{\textbf{sat}}\left ( n \right )}{C_{0}}\textrm{d}n
    \label{eq:DELTA}
\end{equation}
where, $C(n)/C_0$ is the value of the breakthrough curve, $N$ is the total number of pore volumes, $C_{\mathrm{wd}}(n)/C_0$ and $C_{\mathrm{sat}}(n)/C_0$ are the values of breakthrough curves under wetting-drying cycles and saturated conditions, respectively. The outcomes are presented in Figure \ref{fig:EffectofBCandN0}, which also includes the results for samples with initial porosities of 0.4 and 0.5. It is observed that with different initial porosities and biochar contents, the wetting-drying cycle can result in more MPs in the effluent, ranging from 0.99\% to 8.74\% (Figure \ref{fig:EffectofBCandN0}a). From Figure \ref{fig:EffectofBCandN0} (b), the influence of the wetting-drying cycle decreases with more biochar content. By fitting the results using a negative linear function, the gradients of the fitting lines increase with a smaller initial porosity. A lower $n_0$ results in less difference because of wetting-drying cycles. However, the variation because of different initial porosities is insignificant when the biochar content reaches 15\%. \par
\begin{figure}[htbp]
    \centering
    \includegraphics[width=1.0\linewidth]{../figures/Ch5/ch5Figures/EffectofBCandN0.jpg}
    \caption{(a) Cumulative proportion of MP particles in the effluent and (b) increased MPs in effluent caused by wetting-drying cycles}
    \label{fig:EffectofBCandN0} % Fig.7
\end{figure}
To describe the fluctuation when the samples were dried and rewetted, the gradient of the breakthrough curves was calculated at each increment of pore volume (Figure \ref{fig:EffectofWDC}). When the biochar content exceeded 10\%, the gradients under these two hydraulic conditions performed similarly. Only the initial gradient between 1.5 and 2.5 pore volumes is more significant than that under saturated conditions. In contrast, for soil mixtures with less biochar content (0, 5\%), see Figure \ref{fig:EffectofWDC} (a) and (b), the gradient of breakthrough curves under two hydraulic conditions is different. Significant breaks in the gradient evolution occur when the samples are dried and rewetted. The change gradually remains stable after 2 to 3 wetting-drying cycles. This indicates that the wetting-drying cycle can induce new adsorption-desorption equilibria between the adsorbate and adsorbent, which is negative for MP removal. \par
\begin{figure}[htbp]
    \centering
    \includegraphics[width=1.0\linewidth]{../figures/Ch5/ch5Figures/EffectofWDC.jpg}
    \caption{Gradient of breakthrough curves evaluated at biochar contents of (a) 0\%, (b) 5\%, (c) 10\%, and (d) 15\%}
    \label{fig:EffectofWDC} % Fig.8
\end{figure}

\subsection{Effect of physical properties: biochar content and initial porosity}
Figure \ref{fig:ContourofBCandN0} presents the number of MPs remaining in each soil-biochar sample with different initial porosities, biochar contents, and hydraulic conditions, which influence the sample’s adsorption performance. With the increase of biochar content from 0 to 15\%, the amounts of retained MPs increase by twofold or even more than threefold. Comparing the difference between 15\% biochar content and pure soil, the positive function of additional biochar is more significant when the samples have lower porosity and undergo wetting-drying cycles. \par
When the porosity is reduced from 0.50 to 0.30, the quantity of retained MPs increases by 17.03\% in the pure sand sample (Figure \ref{fig:ContourofBCandN0}a). Considering the wetting-drying cycle, only 14.37\% of MP particles were detected in the loose soil column ($n_0$ = 0.50), which increased by 22.02\% when the porosity reached 0.30. However, when the biochar content is 15\%, the increased rate is only 7.41\% and 8.3\% for saturated and wetting-drying conditions, respectively. An empirical model proposed by DeNovio et al. (2004) was used to analyse the effect of porosity, which demonstrated that the most efficient filtration occurs when the size of the MP particles (\qty{1}{\micro\meter}) exceeds 5-10\% of the sample's average grain size (D50). In this study, the diameter of MP particles accounts for only 0.36\% of D50 (0.38 mm), below the particle size required for effective filtration. This suggests that the \qty{1}{\micro\meter} MPs can easily escape from the soil-biochar samples. \par
\begin{figure}[htbp]
    \centering
    \includegraphics[width=1.0\linewidth]{../figures/Ch5/ch5Figures/ContourofBCandN0.jpg}
    \caption{Proportion of trapped MPs in the biochar-sand sample under (a) constant saturated and (b) wetting-drying cycles conditions}
    \label{fig:ContourofBCandN0} % Fig.9
\end{figure}

\subsection{Effect of mixing format}
Three columns (LL05, LL10, LL15) were designed as a sandwich-layered structure consisting of biochar and soil. The column test was only performed under constant saturated conditions for the layered biochar-soil sample, where a pure biochar layer was positioned in the middle. The column without biochar (ML0) was designed as a control group. From Figure \ref{fig:EffectofFormat}, the biochar-soil sample in a sandwiched structure shows a higher retention capacity than the mixed samples. As the biochar content increased from 5\% to 15\%, more MP spheres were retained in the sandwich-format column. In comparison to the scenario with pure sand (ML0), only 5\% biochar content (LL05) in a stratified form resulted in a doubling of the adsorption capacity (45.45\%). This performance is nearly equivalent to that observed in a directly mixed sample (46.24\%) with 10\% biochar content. Additionally, the overall permeability of multi-layered structures is primarily determined by the permeability of the layer with the lowest permeability. Therefore, with the same degree of compaction, the permeability of the layered biochar-soil column is lower than that under mixed form. Thus, the layered structure prompted MP suspension to accumulate directly at the biochar layer, providing enough sites and time for adsorption. \par
\begin{figure}[htbp]
    \centering
    \includegraphics[width=0.65\linewidth]{../figures/Ch5/ch5Figures/EffectofFormat.jpg}
    \caption{Proportion of trapped MPs in the biochar-sand sample in different soil-biochar formats}
    \label{fig:EffectofFormat} % Fig.10
\end{figure}
Figure \ref{fig:MPinLayers} presents the retained MPs in the soil column, the distribution of MPs exhibited a declining trend with depth. O'Connor et al. (2019) demonstrated that a higher concentration of MPs led to considerable adsorption, resulting in a deeper penetration depth. In this study, the concentration of the MP suspension decreased with depth during the adsorption process. Therefore, the probability of contact between the MP particles and the adsorption sites decreases along the depth. MPs accumulated in the biochar part for the layered columns (Figure \ref{fig:MPinLayers}b). The proportion of retained MPs in the biochar layer increased by 47.2\%. This trend is consistent with the previous study by Hsieh et al. (2022). Fig. 11(c) presents the retained MP particles in the unit mass of the samples. The adsorption capacity of the biochar layer is nearly tenfold that of sand. It should be taken for granted that the MP amount in unit biochar decreases with more biochar mixed. However, the adsorption capability in the unit mass of the biochar decreased by 36.81\% when the biochar content changed from 5\% to 10\%, whereas it decreased by only 12.58\% when the biochar content reached 15\%. Therefore, the usage efficiency of biochar is not linearly related to the biochar content. \par
\begin{figure}[htbp]
    \centering
    \includegraphics[width=1.0\linewidth]{../figures/Ch5/ch5Figures/MPinLayers.jpg}
    \caption{Vertical distribution of MPs in the biochar-sand sample under (a) mixed and (b) layered formats, and (c) MPs contents in the unit mass of sample}
    \label{fig:MPinLayers} % Fig.11
\end{figure}

\begin{figure}[htbp]
    \centering
    \includegraphics[width=1.0\linewidth]{../figures/Ch5/ch5Figures/SWRC.jpg}
    \caption{Water retention curves during drying (Initial porosity is 0.5)}
    \label{fig:SWRC} % Fig.12
\end{figure}

\subsection{Water retention characteristic of the soil-biochar sample}
The wetting-drying cycles causes micro-cracks within the soil-biochar samples (Stirling et al., 2021; Li et al., 2023) due to the tensile stress among particles during desiccation, which subsequently causes the MP particles to desorb from adsorption sites. As shown in Figure \ref{fig:FTIRbiochar}, the water retention curves of the samples reduce with the number of wetting-drying cycles. When the biochar content is 0\%, 5\%, 10\% and 15\%, the drying curves approach stability after the third, second, first and first cycle(s). This is consistent with the number of discontinuities in Figure \ref{fig:EffectofWDC}. The water retention ability of the soil-biochar samples increases with higher biochar content. To analyse biochar’s hydrophilic properties, FTIR analysis was employed to analyse the biochar's functional groups (see Figure \ref{fig:FTIRbiochar}). The results demonstrate a broad peak between 3300 and 3100 cm-1, indicative of hydrophilic groups (-NH and -OH) and peaks at 3026 and 2918 cm-1, corresponding to the -CH group vibrations. Therefore, a stable water film can be generated at the surface of biochar particles (Pardo et al., 2019), enhancing the sample’s water retention capability. \par  % 这部分要详细修改，其中关于biochar的ftir结果部分
\begin{figure}[htbp]
    \centering
    \includegraphics[width=0.65\linewidth]{../figures/Ch5/ch5Figures/FTIRbiochar.jpg}
    \caption{FTIR spectra of biochar}
    \label{fig:FTIRbiochar} % Fig.13
\end{figure}


\section{Discussion} \label{sec:5.4}
From the results mentioned above, biochar effectively removes MPs from water and significantly impedes their penetration through the soil medium under cyclic wetting-drying conditions. After the leaching test, an SEM test was performed to demonstrate how MPs were retained in biochar. Figure \ref{fig:SEMmpsinside} (a) indicates that numerous MPs were investigated to accumulate in the tiny crevices and alveolate chambers of biochar particles. These pores have a diameter of around \qtyrange{1}{8}{\micro\meter}, through which MP particles can easily enter, assemble, and become stuck, trapped, and entangled in biochar particles (Wang et al., 2020). Biochar’s porous structure also provides more intricate pathways for seepage (Johnson, 2020). As such, more MPs attach to the biochar’s surface. Moreover, as demonstrated in this study, biochar-amended soil has a better water retention capacity. The internal water meniscus can apply significant surface tension between the soil and biochar particles under unsaturated conditions. This indicates more stable particle alignment in the presence of moisture changes. According to the laws of thermodynamics, pore water preferentially fills smaller pores (Mualem, 1976). Therefore, as shown in Figure \ref{fig:SEMmpsinside} (b), MPs are more prone to assemble at micropores, which essentially reflects the flow tendency of MP suspensions during the drying process (in the unsaturated state). \par
\begin{figure}[htbp]
    \centering
    \includegraphics[width=1.0\linewidth]{../figures/Ch5/ch5Figures/SEMmpsinside.jpg}
    \caption{SEM images of the sample after leaching test: (a) MPs distribution in the biochar-soil sample and (b) MPs attachment on biochar particles}
    \label{fig:SEMmpsinside} % Fig.14
\end{figure}
Figure \ref{fig:Cracks} presents the visual cracks investigated at the sample surface. After five wetting-drying cycles, significant desiccation cracks appeared on the lateral surfaces of the soil, oriented both horizontally and vertically. As the top of the sample is directly exposed to the external environment and serves as the primary area for drying and evaporation, the cracks are most pronounced. The vertical cracks create channels for water flow, thus facilitating the vertical migration of microplastics within the medium (Bläsing and Amelung, 2018). Desiccation cracks in soil media are induced by the internal matrix suction, which manifests as drying shrinkage tensile stress during the drying process. Since soil is considered a granular material with very low tensile strength, cracks will occur when the tensile stress exceeds the soil's tensile capacity (Tang et al., 2011). Correspondingly, biochar, with its unique porous morphology and chemical properties, effectively inhibits desiccation cracks, thereby enhancing the medium's capacity to immobilise microplastics. This is attributed to biochar’s function in amending the soil’s water retention behaviour, thus mitigating crack generation (Lu et al., 2021). The mechanisms are depicted in Figure \ref{fig:Mechanicsm}. On the one hand, the hydroxyl functional groups at the biochar's surface can generate water molecule-enveloped layers, especially under a fully saturated condition (Pardo et al., 2019). The biochar particles release water slowly when the samples gradually become unsaturated. The microscopic pores at the biochar surface contribute to capillary suction. Therefore, during drying, the biochar-amend samples have higher water content and lower tensile stress (Lu et al., 2021). As a result, the propagation of desiccation cracks is reduced, and more MP particles can be trapped. On the other hand, the biochar's porous morphology can contribute to the MP particles’ stability within samples. In this study, the retained MPs can be divided into two categories: sediment MPs among particles (Dong et al., 2018; Rillig et al., 2019) and trapped MPs within biochar pores. As a granularly porous medium, external forces and hydraulic disturbances can result in the dislocation of particles and initiate the desorption of sediment MP particles (Tang et al., 2021). In contrast, the MPs within biochar pores are relatively stable, like being "protected" by biochar particles against mechanical and hydraulic disturbance. Furthermore, the angular shape of biochar particles allows them to easily interlock and entangle with one another, which restrains the relative rotation and displacement of the particles. Moreover, biochar particles occupy the space between soil particles that would otherwise be occupied by shrinkage (Lu et al., 2021). Therefore, the granular structure of soil-biochar samples is more stable and thus performs better crack resistance. With the decrease of initial porosity, the interlock among particles becomes firmer. The MP particles are stably retained during wetting-drying cycles, especially in denser samples. \par
The adsorption efficiency of MP particles is sensitive to the seepage velocity of MP suspension (Tong et al., 2020), which is determined by the soil-biochar sample's permeability. In a loose sample, the attached MP particles are easily eroded. However, the flow rate is lower in a denser sample, and flow paths are more complicated due to the presence of more microscopic pores (Tong et al., 2008; Zhang et al., 2021a). The MP particles are less prone to detachment by slow fluid in a dense medium. Besides, with the rise of biochar in unit volume, the probability that MPs encounter biochar particles increases (Hsieh et al., 2022). \par
\begin{figure}[htbp]
    \centering
    \includegraphics[width=0.6\linewidth]{../figures/Ch5/ch5Figures/Cracks.jpg}
    \caption{Visual cracks on the sample after wetting-drying cycles ($n_0$= 0.5, biochar content: 0)}
    \label{fig:Cracks} % Fig.15
\end{figure}

\begin{figure}[htbp]
    \centering
    \includegraphics[width=1.0\linewidth]{../figures/Ch5/ch5Figures/Mechanicsm.jpg}
    \caption{Mechanism of improving MPs adsorption of sand under hydraulic cycles}
    \label{fig:Mechanicsm} % Fig.16
\end{figure}
This chapter demonstrates that the layered structure of soil-biochar samples offers significant advantages in the adsorption of MP particles under laboratory conditions, as tested on an acid-base soil. This indicates biochar’s potential engineering applications, such as in sewage treatment plants or along highway embankments, where the interlayered biochar can be constructed as a shield layer to prevent MP particles from vertically penetrating (Wang et al., 2020; Ahmad et al., 2023). Besides, most MP particles aggregate within the biochar layer. It is straightforward to recollect and replace the biochar layer when the quantity of MP particles approaches its retention capacity (Hsieh et al., 2022). Therefore, biochar demonstrates significant economic value for engineering practice. \par

\section{Summary and conclusion} \label{sec:5.5} % 这部分内容可以再增加一些
This study investigates the effects of wetting-drying cycles, biochar content, initial porosity and soil-biochar structure on the migration and retention behaviours of PS microplastics in soil-biochar porous medium. The main conclusions drawn from this work are as follows: \par
(1)	The microcracks induced by wetting-drying cycles deteriorate the medium’s retention capability. Therefore, it is an indispensable ambient factor in evaluating the medium's removal efficiency, reflecting long-term reliability under cyclic hydraulic loadings. \par
(2)	As a modification additive, biochar remarkably enhances the removal efficiency of micrometre-scaled MP spheres. The innovation of this study lies in revealing biochar's potential to enhance the soil's porous medium’s adsorption resilience to MPs under cyclic hydraulic loading, thereby improving resistance to desiccation cracks. \par
(3)	Lower porosity diminishes the permeability and provides more available adsorption sites. The slow fluid velocity mitigates the desorption rate. Besides, a compact structure denotes intensified inter-particle embedment and interlocking. This constrains the cracking space to improve the stability of the retention capacity under wetting-drying cycles. \par
(4)	A sandwich-layered format of biochar and sand was demonstrated to exhibit superior removal efficiency, suggesting benefits for simplified engineering construction, despite a decline in permeability. \par
